% Generated draft: requires mathematical and source review.
\documentclass[11pt]{article}
\usepackage[margin=25mm]{geometry}
\usepackage{fontspec}
\setmainfont{Latin Modern Roman}
\usepackage{amsmath,amssymb,mathtools,unicode-math}
\setmathfont{Latin Modern Math}
\usepackage{enumitem,needspace,xurl,hyperref}
\setlength{\parindent}{0pt}
\setlength{\parskip}{5pt}
\setlength{\emergencystretch}{3em}
\newcommand{\N}{\mathbb{N}}
\newcommand{\Z}{\mathbb{Z}}
\newcommand{\Q}{\mathbb{Q}}
\newcommand{\R}{\mathbb{R}}
\newcommand{\C}{\mathbb{C}}
\newcommand{\F}{\mathbb{F}}
\newcommand{\A}{\mathbb{A}}
\newcommand{\PP}{\mathbb{P}}
\newcommand{\E}{\mathbb{E}}
\newcommand{\eps}{\varepsilon}
\newcommand{\dd}{\,\mathrm{d}}
\begin{document}
\section*{UnsolvedMath ID: 701533}
\textbf{Source problem code:} catalogue entry
\subsection{Definitions and mathematical statement}
Let \(p\) be a prime, and let \(\F_p\) be the field with \(p\) elements. Write
\[
M=\left\{(a_i)_{i\geq 1}:a_i\in\Z_{\geq 0},\ a_i=0\text{ for all but finitely many }i\right\},
\qquad x^a=\prod_{i\geq 1}x_i^{a_i}.
\]
Following the explicit description in the supplied source, define the commutative formal power series algebra
\[
R=\left\{\sum_{a\in M}c_a x^a:c_a\in\F_p\right\},
\]
with no restriction on the number of nonzero coefficients. Addition is coefficientwise, and multiplication is given by
\[
\left(\sum_{a\in M}c_a x^a\right)\left(\sum_{b\in M}d_b x^b\right)
=\sum_{e\in M}\left(\sum_{a+b=e}c_a d_b\right)x^e.
\]
For each fixed \(e\), the inner sum is finite. Let
\[
I=\left\{\sum_{a\in M}c_a x^a\in R:c_0=0\right\}
\]
be the maximal ideal of series with zero constant coefficient. Define the abstract ideal generated by products of two elements of \(I\) by
\[
J=I^2=\left\{\sum_{k=1}^{m}f_k g_k:m\in\Z_{\geq 0},\ f_k,g_k\in I\right\}.
\]
Here the sum is finite, and the empty sum is zero; no closure operation is included in the definition of \(J\).

\textbf{Question (John S. Wilson, Kourovka Notebook 18.105).} For a prime \(p\) and the algebra \(R\) above, does there exist an ideal \(U\) of \(R\) such that
\[
U\subsetneq I
\qquad\text{and}\qquad
U+J=I?
\]
The prime is an arbitrary parameter; the ideal \(U\), if it exists, may depend on \(p\).

The source additionally states that an affirmative answer makes \(\mathrm{SL}_3(R)\) and the kernel of the reduction homomorphism
\[
\mathrm{SL}_3(R)\longrightarrow\mathrm{SL}_3(R/U)
\]
provide an answer to the preceding question. Here \(\mathrm{SL}_3(A)\) denotes the group of \(3\times3\) matrices of determinant one over a commutative ring \(A\). The preceding question is not included in the supplied record, so the intended group-theoretic consequence cannot be stated completely from this evidence.
\subsection{Short English statement}
For the algebra of arbitrary formal series in countably many commuting variables over \(\F_p\), let \(I\) be the ideal of series with zero constant term. Wilson asks whether a proper ideal \(U\subsetneq I\) can satisfy \(U+I^2=I\), where \(I^2\) consists of finite sums of products of elements of \(I\). The source's further application refers to an unavailable preceding question.
\subsection{Sources}
\begin{itemize}
\item Local UnsolvedMath registry, permanent ID \(701533\), exact source problem code \texttt{catalogue entry}, title \emph{Kourovka Notebook Problem 18.105}. The supplied record has upstream revision \texttt{d371ee917ac9647fbb06cd6fc529deec6f3dda82} and records the status as \emph{open}; this is supplied registry metadata, not an independently verified mathematical or literature assessment.
\item The supplied attribution is to John S. Wilson, proposed in 2014. The background identifies \emph{The Kourovka Notebook (Unsolved Problems in Group Theory)}, 18th issue (2014), and the 21st issue edition (2026), \texttt{arXiv:1401.0300v48}, original problem number 18.105. The supplied collection description names E. I. Khukhro and V. D. Mazurov as editors of the Issues 1--21 collection. These bibliographic details are reproduced from the record and have not been independently checked.
\end{itemize}
\end{document}
